\section{Metodologia}
\label{sec:metodos}

Trata-se de uma pesquisa de abordagem quantitativa e, quanto aos objetivos, descritiva, na acepção de \citet{gil2017}: busca descrever as características de uma população (os municípios maranhenses) e estabelecer relações entre variáveis, sem manipulá-las. Quanto aos procedimentos técnicos, classifica-se como pesquisa documental \citep{gil2017}, pois utiliza exclusivamente dados secundários de acesso público, produzidos pelo INEP e pelo Instituto Brasileiro de Geografia e Estatística (IBGE) com outras finalidades e aqui reanalisados. O \textit{pipeline} analítico segue as etapas do processo CRISP-DM (\textit{Cross-Industry Standard Process for Data Mining}) \citep{shearer2000}, adaptado ao contexto de análise de dados educacionais: compreensão do domínio, coleta e preparação dos dados, modelagem e avaliação dos resultados. Dentro desse processo, a análise é conduzida em dois estágios sequenciais, resumidos na Figura~\ref{fig:fluxo}: o primeiro responde à QP1 e o segundo, à QP2.

\begin{figure}[!ht]
    \centering
    \small
    \fbox{\parbox{0.92\linewidth}{\centering
        \textbf{Dados}: IDEB 2025 $+$ Censo Escolar 2025 $+$ Censo Demográfico 2022\\[2pt]
        $\downarrow$ integração por município (n = 216)\\[2pt]
        \textbf{Estágio 1 (QP1)}: \textit{K-Means} sobre renda e pobreza extrema;\\ escolha de $k$ por \textit{Silhouette}, estabilidade e ablação\\[2pt]
        $\downarrow$ \textit{clusters} socioeconômicos\\[2pt]
        \textbf{Estágio 2 (QP2)}: em cada \textit{cluster}, quartil superior $\times$ inferior de IDEB;\\ Mann-Whitney, $\delta$ de Cliff, valor $q$, Spearman e validação preditiva
    }}
    \caption{Desenho da análise em dois estágios. Fonte: elaboração própria.}
    \label{fig:fluxo}
\end{figure}

\subsection{Fontes de dados}

Este estudo integra três bases de dados públicas referentes aos municípios do Maranhão:

\begin{itemize}
    \item \textbf{IDEB EFAI 2025 (INEP).} O IDEB dos anos iniciais do ensino fundamental (EFAI), divulgado pelo \citet{inep2025}, é calculado conforme a metodologia proposta por \citet{fernandes2007}: combina a taxa de aprovação escolar (fluxo) com a proficiência média em Língua Portuguesa e Matemática aferida pelo Saeb. Para este estudo, utilizou-se a edição 2025 da rede municipal, que cobriu 216 dos 217 municípios do estado; Nova Iorque não teve o indicador divulgado nesta edição, porque o INEP não divulgou as notas do município no Saeb.

    \item \textbf{Censo Escolar 2025 (INEP).} Os microdados do Censo Escolar \citep{inep2026microdados} fornecem informações em nível de escola sobre infraestrutura física (internet, biblioteca ou sala de leitura, laboratório de informática, quadra esportiva, abastecimento de água e esgotamento sanitário) e sobre o corpo docente. Para as variáveis de formação docente nos anos iniciais, utilizou-se a Tabela 2.29 da Sinopse Estatística da Educação Básica 2025 \citep{inep2026sinopse}. Ela fornece o número de docentes EFAI por nível de formação com numerador e denominador consistentes, sem a dupla contagem presente nos microdados brutos.

    \item \textbf{Censo Demográfico 2022 (IBGE).} Do Censo Demográfico 2022 do IBGE foram extraídas duas medidas municipais de renda domiciliar \textit{per capita}. A primeira é o valor do rendimento nominal médio mensal domiciliar \textit{per capita} (Sistema IBGE de Recuperação Automática, Sidra, tabela 10295) \citep{ibge2022}; a segunda, a distribuição dos moradores por classes de rendimento domiciliar \textit{per capita} em salários mínimos (Sidra, tabela 10296) \citep{ibge2022classes}. Por ser um recenseamento, e não uma pesquisa amostral, o Censo Demográfico é a única fonte de renda com cobertura de todos os 217 municípios do estado.
\end{itemize}

\subsection{Variáveis de segmentação socioeconômica}

A segmentação dos municípios utiliza duas variáveis, ambas do Censo Demográfico 2022:

\begin{itemize}
    \item \textbf{Renda \textit{per capita}}: rendimento nominal médio mensal domiciliar \textit{per capita}, em reais. Mede o \emph{nível} de recursos econômicos disponíveis no município.
    \item \textbf{Pobreza extrema}: percentual de moradores em domicílios com rendimento \textit{per capita} de até 1/4 do salário mínimo ou sem rendimento. Mede a \emph{concentração} de famílias em privação severa. O corte de 1/4 do salário mínimo é o critério de miserabilidade adotado pela Lei Orgânica da Assistência Social para o Benefício de Prestação Continuada \citep{brasil1993}, o que o torna uma linha de referência institucional reconhecida no Brasil.
\end{itemize}

A escolha dessas duas variáveis apoia-se em três argumentos. O primeiro é teórico: o nível socioeconômico é o preditor mais consistente do desempenho escolar desde \citet{coleman1966}, e a meta-análise de \citet{sirin2005} mostra que essa associação é mais forte quando medida em unidades agregadas, como escolas e distritos, do que no nível individual. Controlar o contexto socioeconômico dos municípios antes de comparar seu desempenho é, portanto, condição para não atribuir a fatores escolares diferenças que decorrem da renda \citep{albernaz2002}. O segundo argumento é que as duas medidas captam aspectos complementares da renda: a média é sensível a rendas altas concentradas em poucos domicílios, enquanto a proporção de moradores em pobreza extrema descreve a base da distribuição. Dois municípios com renda média semelhante podem ter proporções bastante diferentes de famílias em privação severa. O terceiro argumento é de consistência temporal: ambas vêm do mesmo recenseamento, de 2022, o mais recente com desagregação municipal completa.

As duas variáveis são correlacionadas (coeficiente de Pearson de $-0{,}87$ entre os 216 municípios analisados), o que é esperado, dado que ambas derivam da distribuição de renda. A contribuição de cada uma para a partição é avaliada no Estágio 1 por meio de um teste de ablação, que refaz a clusterização sem cada variável e mede o quanto a partição se altera.

\subsection{Variáveis}

O conjunto completo de variáveis é organizado em quatro grupos, descritos na Tabela~\ref{tab:variaveis}.

\begin{table*}[!ht]
    \caption{Variáveis utilizadas na análise. As variáveis escolares referem-se à rede municipal e são agregadas por município. Fonte dos dados: Censo Demográfico 2022 (IBGE), Censo Escolar 2025 e Sinopse Estatística 2025 (INEP).}
    \label{tab:variaveis}
    \centering
    \begin{tblr}{%
        colspec = {X[1.1,l] X[0.45,l] X[1.6,l]},
        cells   = {font=\small},
        row{1}  = {font=\small\bfseries},
        width   = 0.92\linewidth
        }
        \toprule
        Variável & Fonte & Definição operacional \\
        \midrule
        \SetCell[c=3]{l}{\textit{Segmentação socioeconômica}} & & \\
        Renda \textit{per capita} & Censo 2022 & Rendimento domiciliar \textit{per capita} médio mensal municipal (R\$) \\
        Pobreza extrema & Censo 2022 & \% de moradores com rendimento domiciliar \textit{per capita} de até 1/4 do salário mínimo ou sem rendimento \\
        \midrule
        \SetCell[c=3]{l}{\textit{Infraestrutura escolar}} & & \\
        \% escolas com internet & Censo 2025 & Proporção de escolas municipais com acesso à internet \\
        \% escolas com biblioteca ou sala de leitura & Censo 2025 & Proporção com biblioteca, biblioteca-sala de leitura ou sala de leitura \\
        \% escolas com laboratório de informática & Censo 2025 & Proporção com laboratório de informática \\
        \% escolas com quadra esportiva & Censo 2025 & Proporção com quadra de esportes \\
        \% escolas com água adequada & Censo 2025 & Proporção com água potável ou rede pública \\
        \% escolas com esgoto adequado & Censo 2025 & Proporção com rede pública ou fossa séptica \\
        \midrule
        \SetCell[c=3]{l}{\textit{Qualificação docente}} & & \\
        \% docentes EFAI com ensino superior & Sinopse 2025 & Proporção de docentes dos anos iniciais com graduação \\
        \% docentes EFAI com licenciatura & Sinopse 2025 & Proporção de docentes dos anos iniciais com licenciatura \\
        \midrule
        \SetCell[c=3]{l}{\textit{Fluxo e distorção}} & & \\
        Média de alunos por turma (educação básica) & Censo 2025 & Total de matrículas da educação básica dividido pelo total de turmas da educação básica \\
        Média de alunos por turma (ensino fundamental) & Censo 2025 & Total de matrículas no ensino fundamental dividido pelo total de turmas do ensino fundamental \\
        \textit{Proxy} de distorção idade-série & Censo 2025 & \% de matrículas com 15 anos ou mais na educação básica municipal \\
        \% turmas com correção de fluxo (anos finais) & Censo 2025 & Proporção de turmas de correção de fluxo entre as turmas dos anos finais do ensino fundamental \\
        \bottomrule
    \end{tblr}
\end{table*}

\subsection{Coleta e pré-processamento dos dados}

A coleta e o tratamento dos dados foram automatizados em \textit{scripts} Python, organizados em três etapas (ingestão, transformação e análise), o que permite reproduzir integralmente a análise. Das três fontes, foram mantidas apenas as informações da rede municipal do Maranhão. No IDEB, os resultados não divulgados pelo INEP foram tratados como valores ausentes. No Censo Escolar, as informações de escola, docente, matrícula e turma foram combinadas escola a escola. Para os docentes dos anos iniciais, usou-se exclusivamente a Sinopse Estatística, de modo que o numerador (docentes com ensino superior ou licenciatura) e o denominador (total de docentes dos anos iniciais) viessem da mesma fonte. Os dados de renda e de pobreza extrema foram obtidos das tabelas do Sidra e armazenados localmente, para que a análise não dependa de alterações posteriores nas bases on-line.

As variáveis de infraestrutura foram agregadas ao nível municipal como proporção de escolas com cada atributo presente; as de docentes, como percentual sobre o total de docentes dos anos iniciais do município; e as de fluxo, como médias ou proporções sobre o total de matrículas e turmas municipais. A unidade de análise é, portanto, o município: as associações encontradas referem-se às redes municipais e não podem ser transpostas para escolas, professores ou alunos individualmente \citep{robinson1950}.

A Tabela 2.29 da Sinopse conta os docentes dos anos iniciais de todas as redes do município (municipal, estadual, federal e privada), enquanto o IDEB analisado é o da rede municipal. Nos \textit{clusters} em que se faz a comparação intragrupo, a rede municipal reúne, em média, 93\% (intermediário) e 95\% (vulnerável) dos docentes dos anos iniciais. Para verificar se esse descompasso afeta os resultados, a formação docente também foi calculada apenas na rede municipal, a partir dos microdados; como os microdados não separam a formação docente por etapa, essa medida abrange toda a educação básica municipal. As três bases foram então integradas por município, resultando em 217 municípios com renda, pobreza extrema e Censo Escolar completos, dos quais 216 têm o IDEB divulgado.

Não há valores ausentes nas variáveis de segmentação nem nas variáveis explicativas; municípios sem turmas de correção de fluxo aparecem com proporção zero nessa variável. A variável-alvo (IDEB) tem um valor ausente (Nova Iorque). Como o município não pode ser classificado por desempenho, ele é excluído dos dois estágios, e a análise segue com 216 municípios.

\subsection{Estágio 1: segmentação socioeconômica por \textit{K-Means}}

O \textit{K-Means} \citep{macqueen1967} agrupa os 216 municípios com base na renda \textit{per capita} e na pobreza extrema. Antes da clusterização, as duas variáveis são padronizadas pela mediana e pelo intervalo interquartil. Essa padronização reduz a influência de municípios atípicos, como São Luís e Imperatriz, cuja renda \textit{per capita} muito acima e pobreza extrema muito abaixo dos demais poderiam distorcer a formação dos grupos.

O algoritmo é executado para um número de grupos $k$ entre 2 e 10, com dez inicializações por valor de $k$ e semente fixa, o que torna o resultado reprodutível. O número de \textit{clusters} é escolhido pelo \textit{Silhouette Score} médio \citep{rousseeuw1987}, que mede a separação relativa dos grupos. A escolha é confirmada pela estabilidade da partição, avaliada por \textit{bootstrap} de 25 iterações com 80\% dos dados, na configuração de referência de \citet{benhur2002}, com o \textit{Adjusted Rand Index} (ARI) como métrica de concordância. A curva de inércia (Método do Cotovelo) é apresentada como apoio visual. A contribuição de cada variável é medida por ablação: a clusterização é refeita sem a variável, e o ARI entre a nova partição e a original indica o quanto a partição depende dela.

A visualização dos \textit{clusters} é feita via PCA \citep{abdi2010} com duas componentes, permitindo inspeção gráfica da separação em espaço reduzido.

\subsection{Estágio 2: análise intragrupo}

Dentro de cada \textit{cluster} com amostra suficiente, os municípios são divididos em dois grupos por desempenho: \textit{alto IDEB}, formado pelo quartil superior ($\geq Q_{75}$), e \textit{baixo IDEB}, formado pelo quartil inferior ($\leq Q_{25}$). Os municípios do intervalo central são excluídos da análise comparativa, o que aumenta o contraste entre os grupos e reduz o ruído. Exige-se um mínimo de dez municípios em cada extremo.

Para cada variável, o teste de Mann-Whitney \citep{mannwhitney1947} bicaudal verifica se as distribuições diferem entre os grupos de alto e de baixo IDEB. A magnitude da diferença é medida pelo $\delta$ de Cliff \citep{cliff1993}, que varia de $-1$ a $1$. Ele indica a probabilidade de um município do grupo alto apresentar valor maior que um do grupo baixo, descontada a probabilidade inversa. Seguindo \citet{vargha2000}, valores absolutos a partir de 0,12, 0,28 e 0,42 são lidos como efeitos pequeno, médio e grande. Como são testadas catorze variáveis em cada \textit{cluster}, os valores de $p$ são acompanhados do valor $q$ de \citet{benjamini1995}, que controla a taxa de falsas descobertas entre os testes. São reportadas as variáveis com $p < 0{,}10$, e a interpretação considera conjuntamente o valor $q$, a magnitude do efeito e a consistência do resultado entre os \textit{clusters}. Por essa razão, e pelo tamanho das amostras, o Estágio 2 tem caráter exploratório.

Duas verificações complementam a comparação, além da medida alternativa de formação docente descrita na seção anterior. A primeira, de robustez, calcula a correlação de Spearman entre cada variável e o IDEB usando todos os municípios do \textit{cluster}, sem descartar o intervalo central, o que testa se as associações dependem do corte por quartis. A segunda avalia o poder preditivo conjunto das variáveis, por validação cruzada em cinco partes repetida vinte vezes. Um classificador \textit{Random Forest} \citep{breiman2001}, com 500 árvores e ponderação das classes, tenta separar os grupos de alto e de baixo IDEB; seu desempenho é medido pela área sob a curva característica de operação do receptor (\textit{Receiver Operating Characteristic}, ROC), abreviada como AUC, em que 0,5 corresponde ao acaso. Uma regressão \textit{ridge} \citep{hoerl1970} tenta prever o IDEB contínuo, no estado e em cada \textit{cluster}; seu desempenho é medido pelo coeficiente de determinação ($R^2$) fora da amostra. As análises usam as bibliotecas pandas, SciPy e scikit-learn.

\subsection{Aspectos éticos}

O estudo usa apenas dados secundários de acesso público, agregados por município e sem possibilidade de identificação de pessoas. Por isso, dispensa apreciação pelo sistema de Comitês de Ética em Pesquisa, nos termos do art.~1º, parágrafo único, da Resolução CNS n.º 510/2016 \citep{cns2016resolucao510}.
